\chapter{考察}
\chapoverview{本章では，第5章で示した結果から得られる知見を整理する．まず，U-Netの出現率差が手法ギャップとして解釈できることを述べる．次に，2015--2018年と2019--2024年の期間差から，後ろ向き評価の設計が成立する可能性を考察する．さらに，本研究の限界として，特許データの限界，語抽出の限界，BERT類似度の解釈上の注意を整理する．最後に，補助ケースとしての発明主題レベル分析と今後の展開を述べる．}

\section{U-Netの手法ギャップとしての解釈}
本研究の結果では，医療画像解析分野におけるU-Net出現率は，製造欠陥検出分野に比べて大きく高かった．Publication numberベースでは，医療画像解析分野のU-Net出現率は5.30\%，製造欠陥検出分野は0.59\%であった．Family IDベースでも，医療画像解析分野は6.78\%，製造欠陥検出分野は0.55\%であった．

この結果は，U-Netが医療画像解析分野で相対的に多く出現し，製造欠陥検出分野では相対的に低出現であることを示す．したがって，U-Netは参照分野で高出現，対象分野で低出現という手法ギャップを持つ技術手法として解釈できる．

ただし，この結果は，U-Netが製造欠陥検出分野に技術的に導入可能であることを直接証明するものではない．本研究で示せるのは，U-Netが専門家確認候補として抽出されうるということである．技術適合性，実装可能性，コスト，性能，権利範囲は，別途専門家が確認する必要がある．

\section{後ろ向き評価の意義}
製造欠陥検出分野におけるU-Net関連特許は，2015--2018年には確認されず，2019--2024年に出現した．この結果は，2018年以前のデータからU-Netを候補として抽出できれば，後年に対象分野で出現した既知事例を過去時点で捉えたことになる，という後ろ向き評価の設計につながる．

本研究の評価設計は，将来を確定的に予測するものではない．むしろ，過去時点の情報から専門家が確認すべき候補を抽出し，その候補が後年の対象分野で実際に出現したかを確認する設計である．この設計により，候補抽出枠組みが既知のクロスドメイン応用事例を捉えられるかを検討できる．

\section{課題・解決手段分離の意義}
全文類似ベクトルのみを用いる場合，特許文書中の課題，解決手段，対象物，効果が混在する．この場合，候補がなぜ選ばれたのかを説明しにくい．本研究では，課題文脈と解決手段文脈を分離することで，候補抽出の理由を次のように説明できる．

第1に，医療画像解析と製造欠陥検出は，画像中の対象領域や異常領域を検出・分割するという点で課題構造が近い．第2に，医療画像解析ではU-Netが高出現であり，製造欠陥検出では過去期間に未出現であった．第3に，後年に製造欠陥検出分野でもU-Net関連特許が出現した．この3点により，U-Netは，課題が近い異分野間で対象分野に低出現であった解決手段の候補として説明できる．

\section{本研究の限界}
本研究には，少なくとも以下の限界がある．

第1に，特許データは技術利用の全体を表すものではない．U-Netは論文，オープンソース実装，実務システムで利用されていても，必ずしも特許として出願されるとは限らない．したがって，特許データ上の低出現は，実務上の未利用を意味しない．

第2に，U-Netの抽出は表記ゆれに依存する．本研究では，U-Net，UNet，U Netを用いたが，別表記や関連構造を完全には捉えられない可能性がある．

第3に，BERT類似度は技術導入可能性を示す指標ではない．BERT類似度は文書上の意味的近さを示すものであり，材料適合性，工程条件，計算コスト，データ制約，権利範囲を評価するものではない．

第4に，U-Netという明示的な技術手法名を主ケースとしているため，洗浄剤や製造方法のように明確な手法名を持たない発明主題レベルの候補抽出には，追加の検討が必要である．

\section{補助ケースとしての発明主題レベル分析}
本研究の補助的展開として，花王の洗浄・除去系技術から半導体洗浄への展開を，発明主題レベルの異質解決策型として扱うことができる．このケースでは，U-Netのような単一手法名ではなく，洗浄剤組成物，基板洗浄方法，残渣除去技術などの解決手段クラスタを対象とする．

この補助ケースは，企業の技術資産活用に近い観点を持つ．一方で，卒論本文の主軸にすると，研究範囲が広がりすぎる危険がある．そのため，本論文では，U-Netを主ケースとし，花王ケースは付録に整理する．

\section{今後の課題}
今後の課題は3つである．

第1に，課題文脈と解決手段文脈の抽出精度を評価する必要がある．代表特許だけでなく，より多くのサンプルについて，LLM抽出結果と人手確認結果を比較する必要がある．

第2に，比較対象を増やす必要がある．本研究ではU-Netを主ケースとしたが，YOLO，GAN，Domain Adaptationなどを補助ケースとして追加することで，提案枠組みの一般性を検討できる．

第3に，特許だけでなく論文データやオープンソースデータを組み合わせる必要がある．特にU-Netのようなソフトウェア系技術では，特許だけでは技術利用の全体像を捉えにくい．論文から特許への展開や，論文・特許・実装の時間差を分析することが今後の重要な課題である．
